\begin{landscape}
\begingroup
\fontsize{10}{13}\selectfont
\fancyhead[L]{\small POSTHUMOUS PAPERS. PERTURBATIONS OF PALLAS}
\renewcommand{\headrulewidth}{0pt}
\par\medskip
\pdfbookmark[2]{Posthumous papers VI [5]}{pallas-VI-fragment5}
\begin{center}[5.]\end{center}
\par\noindent\hspace*{\parindent}[We obtain the perturbations of the semimajor axis in exactly the same way as in the first calculation. Our example works out as follows if we set the constant \(=-\tfrac{2}{8}\lambda\dfrac{10000000}{206265}\), that is, if we choose the 7th\ decimal place as the unit:
\SourceLeaf{556}
\par
\begingroup\small
\noindent\begin{minipage}[c]{.43\textwidth}
\begin{tabular}{@{}lr@{}}
\(b\)\enspace .\,.\,.\,.\,.\,.\,.\,.\,.\,.\,.\,.\,.&\(0{,}95337\)\\
\(b'\)\enspace .\,.\,.\,.\,.\,.\,.\,.\,.\,.\,.\,.\,.&\(0{,}52101\)\\
\multicolumn{2}{r}{\(Q=69^0\,43'\,11''\)}\\
\(\sin Q\)\enspace .\,.\,.\,.\,.\,.\,.\,.\,.&\(9{,}97221\)\\
\(q\)\enspace .\,.\,.\,.\,.\,.\,.\,.\,.\,.\,.\,.\,.&\(0{,}98116\)\\
\(2-5\dfrac{n'}{n}\)\enspace .\,.\,.\,.\,.\,.&\(8{,}74430\)\\
\hline
\(*\)\enspace .\,.\,.\,.\,.\,.\,.\,.\,.\,.\,.\,.\,.&\(2{,}23686\)\\
constant\enspace .\,.\,.\,.\,.\,.\,.&\(1{,}14727_n\)\\
\hline
&\(3{,}38413_n\)
\end{tabular}
\end{minipage}\hfill
\begin{minipage}[c]{.55\textwidth}
hence the perturbation term
\[
-2422\cos(5M'-2M+69^0\,43'\,11'').
\]
\end{minipage}\par
\endgroup
\par\noindent\hspace*{\parindent}The complete result is given together with the perturbations of the first part of the epoch in the following note, since the arguments there are the same. In every case the sine of the corresponding argument is to be taken, so that, for\ example, the perturbation term just calculated reads:
\[
+2422\sin(5\text{\mdseries\AstroSymbols ♃}-2\text{\mdseries\AstroSymbols ⚴}+165^0\,35'\,33'').]
\]
\clearpage
\endgroup
\end{landscape}
